\documentclass[../../main.tex]{subfiles}
\begin{document}

\begin{center}
  {\fontsize{48pt}{58pt}\selectfont \textbf{2009}\normalsize\textbf{年度}}
\end{center}

\begin{table}[htb]
  \centering
  \caption{出題テーマと難易度}
  \begin{tabular}{|c|p{9cm}|c|}
    \hline
    \textbf{問題番号} & \textbf{テーマ} & \textbf{難易度} \\
    \hline
    第1問 & 二項係数の最大公約数と整数の合同 & \\
    \hline
    第2問 & 行列の累乗と収束 & \\
    \hline
    第3問 & 色の異なる玉を箱に入れる確率 & \\
    \hline
    第4問 & 円板の回転でできる立体の体積 & \\
    \hline
    第5問 & 指数関数の不等式 & \\
    \hline
    第6問 & 正三角形から動く点の位置 & \\
    \hline
  \end{tabular}
\end{table}

\begin{table}[htb]
  \centering
  \caption{解答}
  \renewcommand{\arraystretch}{1.4}
  \begin{tabular}{|c|c|p{9cm}|}
    \hline
    \textbf{問題番号} & \textbf{小問} & \textbf{答え} \\
    \hline
    第1問 & (1) & $m$が素数なら$d_m=m$（証明問題） \\
    \cline{2-3}
    & (2) & $k^m-k$は$d_m$で割り切れる（証明問題） \\
    \cline{2-3}
    & (3) & $m$が偶数なら$d_m=1$または$2$（証明問題） \\
    \hline
    第2問 & (1)〜(3) & \\
    \hline
    第3問 & (1) & $P_1=\dfrac{3^2\cdot5^2}{4^6}$ \\
    \cline{2-3}
    & (2) & $P_2=\dfrac{15}{4^3}$ \\
    \cline{2-3}
    & (3) & $\dfrac{P_3}{P_1}=\dfrac{63}{16}$ \\
    \hline
    第4問 & (1) & $W(a)=\dfrac{2}{3}\pi$ \\
    \cline{2-3}
    & (2) & $\displaystyle\lim_{a\to\infty}V(a)=\dfrac{2}{3}\pi$ \\
    \hline
    第5問 & (1) & 指定された指数不等式を示す（証明問題） \\
    \cline{2-3}
    & (2) & $0.9999^{101}<0.99<0.9999^{100}$（証明問題） \\
    \hline
    第6問 & (1) & $|\sin\theta|\le\dfrac{1}{1000}$（証明問題） \\
    \cline{2-3}
    & (2) & $\dfrac{\pi}{3}-2\alpha\le\theta_1+\theta_2
    \le\dfrac{\pi}{3}+2\alpha$ \\
    \cline{2-3}
    & (3) & 時刻$T=\dfrac{1000}{\sqrt{3}}$に3点とも中心から距離$3$以内
    （証明問題） \\
    \hline
  \end{tabular}
\end{table}

\paragraph{試験形態}
6問．試験時間は要確認．

\end{document}
